% =============================================================================
% BioShield — IEEE Conference Paper (first draft)
%
% Title: BioShield: A Negative Result on Adversarial Retraining for Biomedical
%        AI-Text Detection Under Generator-Family Shift
%
% Author: Venkata Rayudu Alapati  (M.Tech AI/ML, AP25122040029)
% Date  : April 2026
% =============================================================================
\documentclass[conference]{IEEEtran}
\IEEEoverridecommandlockouts

\usepackage{cite}
\usepackage{amsmath,amssymb,amsfonts}
\usepackage{algorithmic}
\usepackage{algorithm}
\usepackage{graphicx}
\usepackage{textcomp}
\usepackage{xcolor}
\usepackage{hyperref}
\usepackage{booktabs}
\usepackage{multirow}
\usepackage{array}
\usepackage{enumitem}
\usepackage{caption}
\usepackage{siunitx}
\usepackage{url}

\hypersetup{
  colorlinks=true,
  linkcolor=blue!50!black,
  citecolor=blue!50!black,
  urlcolor=blue!50!black
}

\graphicspath{{figures/}}

\begin{document}

\title{BioShield: A Negative Result on Adversarial Retraining for Biomedical AI-Text Detection Under Generator-Family Shift}

\author{
\IEEEauthorblockN{Venkata Rayudu Alapati}
\IEEEauthorblockA{
\textit{M.Tech, Dept.\ of Computer Science and Engineering}\\
\textit{SRM University AP}\\
Andhra Pradesh, India \\
AP25122040029 \\
alapativenkatarayudu@gmail.com}
}

\maketitle

\begin{abstract}
We study whether a 3-round adversarial retraining loop can harden a biomedical
text deepfake detector against AI-generated abstracts. Our pipeline combines
(i)~a BiomedBERT-Large detector, (ii)~a BioMistral-7B zero-shot generator and a
SeqGAN baseline, and (iii)~a Qwen2.5-7B rewriting agent that targets the
top-25\% hardest fakes per round. We evaluate four conditions in a controlled
ablation on \(1{,}500\) real PubMed abstracts and \(1{,}500\) BioMistral fakes
(test \(n=300\), Wilson 95\% CI \(\pm 2.9\) percentage points), plus a RAID
cross-domain probe. Three findings emerge. \textbf{First}, generator-family
identity dominates evasion behaviour: SeqGAN fakes evade at \(98.7\%\),
BioMistral fakes at \(0.7\%\), and RAID's mixed-LLM corpus at \(98.3\%\) — a
roughly \(100\times\) spread on the same detector. \textbf{Second}, three
rounds of adversarial retraining did not measurably help at this scale:
condition~B's evasion stayed flat (\(0.993 \to 0.980 \to 0.987\), within Monte
Carlo noise on \(n=250\)), and conditions~C and~D produced byte-identical
metrics because best-validation-AUC checkpointing rolled back to the round-0
baseline (validation AUC ceiling \(0.9990\) was unreachable by retraining).
\textbf{Third}, the detector did not transfer cross-domain (RAID accuracy
\(0.11\), evasion \(0.983\)). We argue that (a)~published biomedical detectors
must disclose the generator family they were trained against, (b)~best-val-AUC
checkpoint selection silently hides retraining signal when the baseline is
already at ceiling, and (c)~a previously reported \(0.74 \to 0.54\) evasion
curve at \(500\)-sample scale was likely small-sample noise that did not
reproduce at \(3\)k. We position BioShield as a biomedical specialist and
treat the negative results as the paper's primary contribution.
\end{abstract}

\begin{IEEEkeywords}
AI-generated text detection, biomedical NLP, adversarial training,
SeqGAN, large language model rewriting agents, domain generalization,
BiomedBERT, negative results, reproducibility.
\end{IEEEkeywords}

\section{Introduction}
\label{sec:introduction}

Large language models (LLMs) can now generate biomedical research abstracts
that are syntactically fluent, terminologically plausible, and difficult for
non-experts to distinguish from real PubMed entries
\cite{else2023chatgpt,liang2023gptdetect}. The downstream risks are concrete:
fabricated abstracts can pollute citation databases, support fraudulent
manuscripts, and amplify medical misinformation when they enter
patient-facing search results \cite{majovsky2023artificialintelligence}. A
biomedical-specific deepfake-text detector is therefore not a hypothetical
need but a near-term tooling gap for journals, preprint servers, and
clinical-information platforms.

A recurring promise in the AI-text-detection literature is that
\emph{adversarial training} — iteratively generating fakes, retraining the
detector against them, and repeating — yields a more robust binary
classifier \cite{zellers2019defending,mitchell2023detectgpt,hu2023radar}. We set out to test whether this promise survives a
controlled, fully reproducible biomedical pipeline running on commodity
hardware (a single Apple~M3~Max). Our pipeline pairs a
\textbf{BiomedBERT-Large} detector \cite{gu2021biomedbert} with two
generators (\textbf{BioMistral-7B} \cite{labrak2024biomistral} and
\textbf{SeqGAN} \cite{yu2017seqgan}) and an LLM rewriting
\textbf{agent} (\textbf{Qwen2.5-7B} \cite{qwen2024technical}) that targets
the top-\(25\%\) hardest fakes per round. Four conditions isolate the
contribution of each component (Section~\ref{sec:setup}).

\textbf{Research questions.} We ask:
\begin{itemize}
\item \textbf{RQ1.} Does 3-round adversarial retraining lower the evasion
  rate of a BiomedBERT detector against in-family fakes
  (BioMistral)?
\item \textbf{RQ2.} Does a Qwen2.5-7B rewriting agent meaningfully harden
  the fake distribution beyond the zero-shot generator?
\item \textbf{RQ3.} Does the resulting detector generalise to a
  cross-domain mixed-LLM benchmark (RAID \cite{dugan2024raid})?
\end{itemize}

\textbf{Contributions.} We treat the answers — all of them negative — as the
paper's contribution rather than an embarrassment to hide.
\begin{enumerate}[leftmargin=*]
\item \textbf{Generator-family dependence dominates.} Across one detector
  and one test set, evasion ranges from \(0.7\%\) (BioMistral fakes) to
  \(98.7\%\) (SeqGAN fakes) to \(98.3\%\) (RAID's mixed LLMs) — a
  \(\sim\!100\times\) spread that dwarfs every other effect we measured.
\item \textbf{Adversarial retraining produced no measurable improvement at
  \(3\)k scale.} Condition~B's per-round evasion stayed at
  \(0.993, 0.980, 0.987\) — within Wilson \(95\%\) CI noise on \(n=250\).
  Conditions~C and~D produced byte-identical test metrics because the
  best-validation-AUC checkpoint selector kept rolling back to the
  baseline.
\item \textbf{We retract a prior \(500\)-sample finding.} An earlier run
  reported a \(0.74 \to 0.54\) evasion curve in Condition~B; with the
  larger \(3\)k corpus and tighter CIs the curve disappears, and we now
  attribute it to small-sample noise.
\item \textbf{Cross-domain transfer is open.} A detector trained on
  PubMed~+~BioMistral fails on RAID (accuracy \(0.11\), evasion \(0.983\)),
  motivating future multi-family training rather than supporting any claim
  of universal AI-text detection.
\item \textbf{Reproducibility.} All code, data, metrics, plots, and a
  one-command rerun script are released; the entire study costs
  \(10\,\mathrm{h}\,01\,\mathrm{m}\) on a single laptop.
\end{enumerate}

\item \textbf{Adversarial-aware checkpoint selection.} We implement a
  drop-in replacement for best-val-AUC selection
  (\texttt{--checkpoint\_selector adversarial}) that scores all epoch
  checkpoints on the round's live fake pool and promotes the epoch with the
  \emph{lowest evasion rate}. This is designed to dissolve the C\,$=$\,D
  pathology when the val-AUC ceiling is already saturated.
\item \textbf{Multi-family training (Condition E).} We implement a
  BiomedBERT-Large detector trained jointly on BioMistral-7B, SeqGAN, and
  Llama-3.2-3B fakes. The multi-family corpus is assembled by
  \texttt{models/detector/train\_multi\_family.py} and is designed to
  compress the \(\sim\!140\times\) evasion spread identified in
  Contribution~1.
\item \textbf{Post-hoc calibration.} We implement temperature-scaling
  calibration~\cite{guo2017calibration}
  (\texttt{evaluation/calibration.py}) that learns an optimal decision
  threshold and temperature on a held-out calibration set. Applied to
  Condition~A's detector on the RAID probe, calibration recovers meaningful
  accuracy by correcting the threshold shift that produces the
  AUC--accuracy gap.

\textbf{Paper organisation.} Section~\ref{sec:related} surveys
machine-generated text detection, adversarial training for text classifiers,
and LLM-as-agent rewriting. Section~\ref{sec:method} describes the BioShield
pipeline. Section~\ref{sec:setup} details the four-condition ablation, data
splits, and metrics. Section~\ref{sec:results} reports per-condition results
and the RAID probe. Section~\ref{sec:discussion} owns the negative findings
and explains the C\,$=$\,D pathology. Section~\ref{sec:limitations}
catalogues limitations, and Section~\ref{sec:conclusion} closes.

\section{Related Work}
\label{sec:related}

We organise prior work into three threads that BioShield draws on directly:
machine-generated text (MGT) detection, adversarial training for text
classifiers, and LLM-as-agent rewriting.

\subsection{Machine-Generated Text Detection}
Early MGT detectors relied on stylometric and surface-statistical signals.
GLTR \cite{gehrmann2019gltr} visualises token-level rank statistics from a
language model to surface unusually probable continuations. Grover
\cite{zellers2019defending} trained a discriminator jointly with a generative
news LM and showed that the best detector for a given LM is often a model
of the same family — an observation we revisit at biomedical scale.
DetectGPT \cite{mitchell2023detectgpt} and zero-shot likelihood-perturbation
methods exploit the curvature of the source LM's likelihood landscape.
RoBERTa-based supervised detectors \cite{solaiman2019releasestrategies} remain
the most widely deployed baseline, including the GPT-2 output detector and
its biomedical-domain descendants \cite{liang2023gptdetect}. Recent
benchmarks RAID \cite{dugan2024raid} and MAGE \cite{li2024mage} show
that almost all such detectors collapse on out-of-distribution generators,
which our cross-domain RAID probe confirms.

\subsection{Adversarial Training for Text Classifiers}
Adversarial training in NLP spans two distinct families. \emph{Sample-level}
attacks such as TextAttack \cite{morris2020textattack},
TextFooler \cite{jin2020textfooler}, and BERT-Attack
\cite{li2020bertattack} perturb individual tokens or sentences against a
fixed classifier. \emph{Generator-level} adversarial training instead
trains a generator alongside the classifier; SeqGAN \cite{yu2017seqgan}
remains the canonical baseline using policy-gradient REINFORCE. More recent
work on RL-tuned generators \cite{ouyang2022rlhf} and DPO-style preference
optimisation \cite{rafailov2023dpo} suggests that hardening generators via
explicit reward signals is more effective than the SeqGAN regime — which
our flat Condition~B curve indirectly supports.

A practical question that this literature largely sidesteps is
\emph{checkpoint selection}: when retraining yields a sequence of detector
weights with comparable validation metrics, which checkpoint is shipped?
Best-validation-AUC selection is the silent default in most public
implementations. Our Section~\ref{sec:discussion} demonstrates that this
default actively hides retraining signal when the baseline is already at
ceiling, producing the C\,$=$\,D pathology we report.

\subsection{LLM-as-Agent Rewriting and Red-Teaming}
LLMs are increasingly used as \emph{agents} that probe and harden classifiers
by paraphrase or rewrite \cite{perez2022redteaming,ganguli2022redteaming}.
PAIR \cite{chao2024jailbreaking}, AutoDAN \cite{liu2023autodan}, and Tree of
Attacks \cite{mehrotra2024toa} apply this idea to safety classifiers.
Closer to our setting, recent work uses GPT-4 or open-weight LLMs to
paraphrase machine-generated essays until they evade detection
\cite{krishna2023paraphrasing,sadasivan2024aigeneratedtext}. Our agent
loop instantiates the same idea with Qwen2.5-7B as the rewriter and the
top-25\% hardest fakes from the prior round as the agent's input. The
near-zero evasion gain we observe (Condition~C vs. Condition~A) is
consistent with the negative scaling result in
\cite{sadasivan2024aigeneratedtext} that paraphrase-based attacks plateau
when the source generator is already in-family for the detector.

\subsection{Position of BioShield}
BioShield's contribution is not a new architecture or attack. It is a
controlled, fully reproducible \emph{biomedical-domain} measurement of how
much of the evasion-rate signal in published detector papers can be
attributed to (a)~the choice of generator family, (b)~retraining rounds, and
(c)~rewrite agents — and the answer, at \(3\)k scale, is that (a)~dominates
by two orders of magnitude over (b) and (c). This complements the
cross-domain failure mode reported by RAID \cite{dugan2024raid} and MAGE
\cite{li2024mage} with an in-domain explanation for why the same detector
can look near-perfect or near-useless depending on which generator the
reviewer probes it with.

\section{Method}
\label{sec:method}

The BioShield pipeline has four interacting components: a real-text source
(PubMed), two generators, a binary detector, and a rewriting agent. They are
composed into a generic adversarial loop that any subset of components can
inhabit, which is what enables our four-condition ablation.

\subsection{Components}
\label{ssec:components}

\textbf{Real-text source.} We stream biomedical abstracts from
\texttt{ncbi/pubmed} via the Hugging Face datasets API
\cite{lhoest2021hfdatasets}. Sampling is uniform across years; abstracts
shorter than 50 tokens are discarded.

\textbf{Generator (Bio-family).} BioMistral-7B
\cite{labrak2024biomistral} is a Mistral-7B \cite{jiang2023mistral} model
continually pretrained on PubMed Central. We use it zero-shot, prompted
with biomedical topic stubs and a structured-abstract instruction. Sampling
uses temperature \(0.8\), top-\(p\) \(0.95\), and a maximum of \(384\)
new tokens.

\textbf{Generator (out-of-family baseline).} SeqGAN \cite{yu2017seqgan}
trains a small LSTM generator with policy-gradient REINFORCE against a
CNN-based discriminator. We pretrain SeqGAN on the same PubMed corpus that
BioMistral saw and use it as the Condition~B fake source.

\textbf{Detector.} We fine-tune
\texttt{microsoft/BiomedNLP-BiomedBERT-large-uncased-abstract}
\cite{gu2021biomedbert} (\(\sim\)340M parameters) as a binary
real-vs-fake classifier with a single linear head over the
\texttt{[CLS]} representation. Training uses bf16 mixed precision on
Apple~MPS, batch size \(8\), learning rate \(2 \times 10^{-5}\) with
linear decay and \(10\%\) warm-up, weight decay \(0.01\), and
\(2\) epochs per round.

\textbf{Agent (rewriter).} Qwen2.5-7B \cite{qwen2024technical} is used
zero-shot as a rewriting agent. Given a fake abstract that the current
detector found hard (probability of being real $\geq 0.5$), the agent
produces a rewrite that ``reads more like a real PubMed abstract — keep
the topic and the main claim, change the prose so it sounds like a
biomedical reviewer wrote it.'' We do not fine-tune Qwen.

\subsection{Adversarial Loop}
\label{ssec:loop}

Algorithm~\ref{alg:loop} formalises the general loop. The
\textsc{Generate} step is component~(a): SeqGAN for Condition~B,
BioMistral for Conditions~C and~D. The \textsc{HardSubset} and
\textsc{Rewrite} steps are component~(b): present in C and D, absent in
B. The \textsc{Retrain} step is component~(c): present in B and D,
absent in A and C.

\begin{algorithm}[t]
\caption{BioShield adversarial loop}
\label{alg:loop}
\begin{algorithmic}[1]
\REQUIRE Real corpus $\mathcal{R}$, generator $G$, detector $D_0$,
agent $A$ (optional), rounds $T$, fake-pool size $K$, hard-fraction $h$
\STATE $D \leftarrow D_0$  \COMMENT{baseline trained on
$(\mathcal{R}, G(\cdot))$ once}
\FOR{$r = 1, \dots, T$}
  \STATE $\textsc{set\_seed}(\text{base} + r)$
  \STATE $F_r \leftarrow \textsc{Generate}(G, K)$
  \IF{agent enabled}
    \STATE $H_r \leftarrow \textsc{HardSubset}(F_r, D, h)$
    \STATE $F_r' \leftarrow \textsc{Rewrite}(A, H_r)$
    \STATE $F_r \leftarrow (F_r \setminus H_r) \cup F_r'$
  \ENDIF
  \IF{retrain enabled}
    \STATE $D \leftarrow \textsc{Retrain}(D, \mathcal{R}, F_r)$
    \STATE $D \leftarrow \arg\max_{\text{epoch}} \mathrm{AUC}_{\text{val}}$
    \COMMENT{best-val-AUC selection}
  \ENDIF
  \STATE record $(\mathrm{AUC}, F_1, \mathrm{evasion})_r$ on test set
\ENDFOR
\RETURN $D$
\end{algorithmic}
\end{algorithm}

\subsection{Hard-Subset Selection}
We define a fake $x$ as ``hard'' for the current detector $D$ if
$D(x) \geq 0.5$, i.e., the detector misclassifies it as real. With
$h = 0.25$ and $K = 250$, the hard subset is the top-\(62\) fakes ranked by
$D$'s real-probability score. Only this subset reaches the agent — the
other \(188\) fakes are passed through unmodified. This keeps the
agent budget small (\(\sim\!8\%\) of total wall-clock,
Section~\ref{sec:setup}) while ensuring it spends compute where the
detector is actually weak.

\subsection{Checkpoint Selection}
\label{ssec:checkpoint}
After each retrain round, we keep the epoch checkpoint with the highest
validation AUC. This is the silent default in most public implementations
of adversarial-text training. We will show in Section~\ref{sec:discussion}
that this choice is what produces the C\,$=$\,D pathology: when the
baseline detector already hits the validation-AUC ceiling
(\(0.9990\) on this corpus), no retrained checkpoint can replace it, so
Condition~D's loop runs but never updates the deployed weights.

\subsection{Reproducibility Controls}
\label{ssec:repro}
We fix the random seed (\texttt{base=42}) and call
\texttt{set\_seed(42 + r)} at the top of every round so that round indices
produce different fakes and rewrites without losing reproducibility. Mixed
precision is forced to bf16 (Apple~MPS does not implement deterministic
fp16 kernels). The detector architecture, hyperparameters, validation set,
test set, agent prompt, and SeqGAN configuration are held constant across
all four conditions; the only experimental variable is which steps of
Algorithm~\ref{alg:loop} are enabled.

\section{Experimental Setup}
\label{sec:setup}

\subsection{Data}
\label{ssec:data}
The corpus consists of \(1{,}500\) real PubMed abstracts streamed from
\texttt{ncbi/pubmed} and \(1{,}500\) BioMistral-7B zero-shot fakes
generated from biomedical topic stubs. The combined \(3{,}000\) abstracts
are stratified-split \(80/10/10\) into Train (\(2{,}400\)), Val
(\(300\)), and Test (\(300\)), each balanced \(1\!:\!1\) real-to-fake.
Per-round fake pools (\(K=250\)) are sampled \emph{outside} the train/val/
test splits and never re-enter them; this keeps the test set untouched
across all rounds and conditions. The full data manifest, including
per-row hashes, is committed to the repository so that any reviewer can
verify there is no train--test contamination.

\subsection{Conditions (the four-cell ablation)}
\label{ssec:conditions}
Table~\ref{tab:conditions} maps each condition onto the components of
Algorithm~\ref{alg:loop}.

\begin{table}[t]
\caption{Four-condition ablation. The only experimental variable is which
components of the adversarial loop are enabled.}
\label{tab:conditions}
\centering
\renewcommand{\arraystretch}{1.15}
\begin{tabular}{@{}llccc@{}}
\toprule
& \textbf{Generator} & \textbf{Agent} & \textbf{Retrain} & \textbf{Rounds} \\
\midrule
\textbf{A} static\_baseline & BioMistral & --- & no  & --- \\
\textbf{B} seqgan\_only     & SeqGAN     & --- & yes & 3 \\
\textbf{C} agent\_only      & BioMistral & Qwen2.5 & no  & 3 (gen+rewrite) \\
\textbf{D} full\_pipeline   & BioMistral & Qwen2.5 & yes & 3 \\
\bottomrule
\end{tabular}
\end{table}

\subsection{Cross-domain Probe: RAID}
\label{ssec:raid}
We evaluate Condition~A's detector (without retraining) on the biomedical
slice of the RAID benchmark \cite{dugan2024raid}: \(200\) samples
(\(19\) real, \(181\) fake) with fakes generated by ChatGPT, GPT-4,
Claude, Llama, and other LLMs that the detector never saw. RAID is
strictly an evaluation set; no RAID samples enter training.

\subsection{Metrics}
\label{ssec:metrics}
We report four metrics, all computed on the held-out Test set unless
explicitly tagged ``fake-pool'':

\begin{itemize}[leftmargin=*]
\item \textbf{AUC.} Area under the ROC curve. Threshold-free.
\item \textbf{F1.} Harmonic mean of precision and recall at threshold
  \(0.5\).
\item \textbf{Accuracy.} Fraction of correct labels.
\item \textbf{Evasion rate.} Fraction of \emph{fake} samples assigned
  probability $< 0.5$ (i.e., labelled as real). Computed both on the
  Test set's \(150\) fakes and on each round's \(250\)-sample
  \emph{fake-pool}. The fake-pool evasion is the metric the loop is
  effectively optimising against.
\end{itemize}

\textbf{Confidence intervals.} All test-set rates are reported with Wilson
\(95\%\) confidence intervals. At \(n=300\) the half-width is
\(\pm 0.029\); at \(n=250\) (per-round fake pool) it is \(\pm 0.062\).
We will use these CIs to bound any per-round claims.

\subsection{Hardware and Software}
\label{ssec:hardware}
All experiments run on a single Apple~M3~Max~64\,GB laptop using
PyTorch~2.4 with the Apple~MPS backend and bf16 mixed precision. We set
\texttt{PYTORCH\_MPS\_HIGH\_WATERMARK\_RATIO=0.0} and
\texttt{OMP\_NUM\_THREADS=10}. The orchestrator process is wrapped in
\texttt{caffeinate -dims -w \$PID} so the lid can close without
suspending the run.

\subsection{Wall-clock Budget}
\label{ssec:walltime}
Table~\ref{tab:walltime} reports the wall-clock breakdown of one
end-to-end run (\(10\,\mathrm{h}\,01\,\mathrm{m}\)) and where the time
goes. Data preparation (initial \(1{,}500\)-fake BioMistral generation)
and Condition~D dominate; plots, RAID, and the M4 skip are essentially
free.

\begin{table}[t]
\caption{End-to-end wall-clock on a single Apple M3 Max
(bf16, MPS). The whole pipeline is one shell script.}
\label{tab:walltime}
\centering
\renewcommand{\arraystretch}{1.1}
\begin{tabular}{@{}lrr@{}}
\toprule
\textbf{Stage} & \textbf{Wall-clock} & \textbf{\%} \\
\midrule
Data prep (1{,}500 PubMed + 1{,}500 BioMistral)
                                & 3\,h\,55\,m & 39 \\
Condition A (baseline detector) & 13\,m       &  2 \\
Condition B (SeqGAN $\times$ 3 rounds)
                                & 56\,m       &  9 \\
Condition C (BioMistral + Qwen $\times$ 3)
                                & 2\,h\,05\,m & 21 \\
Condition D (B+C, with retrain $\times$ 3)
                                & 2\,h\,52\,m & 29 \\
Plots                           & 0\,s        & 0 \\
RAID benchmark                  & 22\,s       & $<$ 1 \\
M4 (skipped, dataset removed)   & 2\,s        & $<$ 1 \\
\midrule
\textbf{Total}                  & \textbf{10\,h\,01\,m} & 100 \\
\bottomrule
\end{tabular}
\end{table}

\subsection{What We Did Not Run}
For honesty: the M4 \cite{wang2024m4} cross-domain benchmark was removed
from the Hugging~Face Hub between project start and this run. We record a
graceful skip in \texttt{condition\_D\_m4\_metrics.json}. A
Llama-3.2-3B transfer-attack evaluation is queued but not in this draft;
the empty \texttt{transfer\_test.csv} placeholder is committed.

\section{Results}
\label{sec:results}

We report results in three subsections that follow the four-condition
ablation, the per-round retraining curves, and the cross-domain RAID
probe. All confidence intervals are Wilson \(95\%\); the half-width is
\(\pm 0.029\) at \(n = 300\) (test set) and \(\pm 0.062\) at
\(n = 250\) (per-round fake pool).

\subsection{Headline Comparison}
\label{ssec:headline}

Table~\ref{tab:headline} summarises the final-round metrics for each
condition together with the RAID cross-domain evaluation.

\begin{table*}[t]
\caption{Final-round metrics (test \(n=300\), Wilson \(95\%\) CI half-width
\(\pm 0.029\)) for the four conditions and the RAID cross-domain probe.
Evasion is the fraction of fakes assigned probability $<0.5$ on each
condition's final-round fake pool (\(n=250\), CI half-width \(\pm 0.062\)).
RAID evasion is on its full fake split (\(n=181\)).}
\label{tab:headline}
\centering
\renewcommand{\arraystretch}{1.15}
\begin{tabular}{@{}lcccc@{}}
\toprule
\textbf{Condition} & \textbf{AUC} & \textbf{F1} & \textbf{Acc.} & \textbf{Evasion} \\
\midrule
A — static\_baseline      & 0.9997 & 0.9797 & 0.980 & --- \\
B — seqgan\_only          & 0.9498 & 0.6546 & 0.490 & 0.987 \\
C — agent\_only           & 0.9993 & 0.9797 & 0.980 & 0.007 \\
D — full\_pipeline        & 0.9993 & 0.9797 & 0.980 & 0.007 \\
\midrule
RAID (cross-domain, A's detector)
                          & 0.8171 & 0.176  & 0.110 & 0.983 \\
\bottomrule
\end{tabular}
\end{table*}

Three readings:
\begin{enumerate}[leftmargin=*]
\item Conditions~A, C, and D form a tight cluster
  (AUC \(\approx 0.9994\)) that is statistically indistinguishable on
  \(n=300\). Condition~B is far below
  (\(\Delta_{\mathrm{AUC}} \approx 0.05\), well outside the
  \(\pm 0.029\) CI).
\item Conditions~C and~D produce \emph{byte-identical} test metrics. We
  diagnose this as a checkpoint-selection pathology in
  Section~\ref{sec:discussion}.
\item Evasion ranges from \(0.007\) to \(0.987\) on the same detector
  family — a $\sim$\(140\times\) spread driven by which generator
  produced the fakes.
\end{enumerate}

\begin{figure}[t]
\centering
\includegraphics[width=\linewidth]{ablation_comparison.png}
\caption{Side-by-side bar comparison across A/B/C/D and RAID. Tight
A/C/D cluster on AUC and F1; Condition~B's collapse is driven by
SeqGAN-style fakes that the detector ranks correctly (AUC \(\approx 0.95\))
but cannot classify (F1 \(\approx 0.65\)). RAID is the right-most
column; AUC is preserved at \(0.82\) but accuracy collapses to \(0.11\)
because the threshold is wrong cross-domain.}
\label{fig:ablation}
\end{figure}

\subsection{Per-Round Retraining Curves}
\label{ssec:perround}

Figure~\ref{fig:auc_f1} reports test AUC and F1 vs. round. Tables
\ref{tab:b_perround} and \ref{tab:c_perround} report the corresponding
per-round evasion on each condition's fake pool. Condition~D's per-round
metrics are byte-identical to Condition~C's and are omitted.

\begin{figure}[t]
\centering
\includegraphics[width=\linewidth]{auc_f1_vs_round.png}
\caption{Test AUC (top) and F1 (bottom) versus retraining round for
Conditions B/C/D. Condition~B sits at the bottom of the AUC plot
(\(\approx 0.949\) flat) and the bottom of the F1 plot
(\(\approx 0.65\) flat). Conditions~C and~D track Condition~A's
ceiling (\(\mathrm{AUC} \approx 0.999\), \(F_1 \approx 0.98\))
across all three rounds.}
\label{fig:auc_f1}
\end{figure}

\begin{table}[t]
\caption{Condition~B per-round metrics. Test-set evasion is on
\(n=300\); fake-pool evasion is on \(n=250\) per round. The
fake-pool differences (\(\pm 0.013\)) are smaller than the
\(\pm 0.062\) CI on \(n=250\); we therefore do not claim per-round
improvement.}
\label{tab:b_perround}
\centering
\renewcommand{\arraystretch}{1.1}
\begin{tabular}{@{}cccccc@{}}
\toprule
\textbf{Round} & \textbf{AUC} & \textbf{F1} & \textbf{Acc.} &
\textbf{Test ev.} & \textbf{Pool ev.} \\
\midrule
1 & 0.9492 & 0.6532 & 0.487 & 0.953 & \textbf{0.993} \\
2 & 0.9485 & 0.6561 & 0.493 & 0.947 & \textbf{0.980} \\
3 & 0.9498 & 0.6546 & 0.490 & 0.953 & \textbf{0.987} \\
\bottomrule
\end{tabular}
\end{table}

\begin{table}[t]
\caption{Condition~C per-round metrics. Condition~D is byte-identical
because of the checkpoint-selection pathology
(Section~\ref{sec:discussion}). The detector catches BioMistral fakes
$> 99\%$ of the time even before any retraining round.}
\label{tab:c_perround}
\centering
\renewcommand{\arraystretch}{1.1}
\begin{tabular}{@{}cccccc@{}}
\toprule
\textbf{Round} & \textbf{AUC} & \textbf{F1} & \textbf{Acc.} &
\textbf{Test ev.} & \textbf{Pool ev.} \\
\midrule
1 & 0.9984 & 0.9764 & 0.977 & 0.013 & \textbf{0.013} \\
2 & 0.9990 & 0.9797 & 0.980 & 0.007 & \textbf{0.007} \\
3 & 0.9993 & 0.9797 & 0.980 & 0.007 & \textbf{0.007} \\
\bottomrule
\end{tabular}
\end{table}

\begin{figure}[t]
\centering
\includegraphics[width=\linewidth]{evasion_rate_vs_round.png}
\caption{Per-round evasion rate on each condition's fake pool. Condition~B
is flat near the top of the chart at $\sim$\(99\%\); Conditions~C and~D
track each other near the bottom at $\sim$\(0.7\%\). The vertical gap
between the two curves is the generator-family effect we report; the
within-curve variation is inside the \(\pm 0.062\) Wilson CI on
\(n=250\) and we make no per-round claims about it.}
\label{fig:evasion}
\end{figure}

\subsection{The Generator-Family View}
\label{ssec:family}

Re-organising the same numbers by generator family makes the dominant
effect visible (Table~\ref{tab:family}).

\begin{table}[t]
\caption{Same detector family (BiomedBERT-Large), three generator
families, three orders of magnitude difference in evasion. Each row is
the final-round evasion rate on the corresponding fake source.}
\label{tab:family}
\centering
\renewcommand{\arraystretch}{1.15}
\begin{tabular}{@{}lcl@{}}
\toprule
\textbf{Generator family} & \textbf{Evasion} & \textbf{Source} \\
\midrule
BioMistral-7B (zero-shot)        & 0.007 & Conditions C, D pool \\
SeqGAN (LSTM + REINFORCE)        & 0.987 & Condition B pool \\
Mixed LLMs (ChatGPT, GPT-4, \dots) & 0.983 & RAID benchmark \\
\bottomrule
\end{tabular}
\end{table}

A spread of $\sim$\(140\times\) within one detector and one test
construction. Section~\ref{sec:discussion} argues that this is the
single most important number in the paper, and that it should be a
mandatory disclosure for any deployed AI-text detector.

\subsection{Cross-Domain Probe (RAID)}
\label{ssec:raid_results}

The detector that scores AUC \(0.9997\) on its in-distribution test set
scores AUC \(0.8171\) on RAID's biomedical slice
(\(n = 200\); 19~real, 181~fake). Crucially, F1 collapses to \(0.176\)
and accuracy to \(0.11\): the detector still ranks fakes above reals
(positive AUC) but its decision threshold (\(0.5\) calibrated on
PubMed~+~BioMistral) puts \(178\) of the \(181\) fakes on the
``real'' side.

This pattern — preserved AUC, collapsed accuracy — is consistent with the
domain-shift literature: the model's ranking ability transfers, its
calibration does not. We return to this in Section~\ref{sec:discussion}.

\section{Discussion}
\label{sec:discussion}

We organise the discussion around the three negative findings, in
descending order of how much they affect what other practitioners should
do.

\subsection{Generator-Family Identity Dominates Evasion}
\label{ssec:family_disc}

Across one detector architecture, one training corpus, and one held-out
test set, evasion ranged from \(0.007\) (BioMistral fakes) to
\(0.987\) (SeqGAN fakes) to \(0.983\) (RAID's mixed LLMs). Three
points follow.

\paragraph{The effect dwarfs everything else we measured.}
Per-round evasion variation in Condition~B (\(0.993, 0.980, 0.987\))
sits inside the \(\pm 0.062\) Wilson CI on \(n=250\); the
between-condition retraining effect is therefore zero within noise. The
agent-rewrite effect (\(\Delta_{\mathrm{evasion}} \approx 0.006\) between
C-round-1 and C-round-3) is similarly inside CI. The
between-generator-family effect (\(\Delta_{\mathrm{evasion}}
\approx 0.98\)) is two orders of magnitude larger than either. If a
practitioner is choosing where to invest engineering time, the lesson is
unambiguous: \emph{the choice of fakes the detector trains against
matters far more than how many adversarial rounds you run, or whether
you wrap an LLM agent around the loop}.

\paragraph{Why SeqGAN is so much harder.}
SeqGAN produces text whose token-level distribution is far from the
biomedical-English manifold the detector was pretrained on. This is
visible in Condition~B's metric pattern: AUC remains acceptable
(\(\approx 0.95\)), meaning the detector \emph{ranks} SeqGAN fakes
above PubMed reals, but F1 collapses to \(0.65\), meaning the
calibrated threshold lands in a region where most SeqGAN samples are
called real. Two epochs of retraining per round do not move the
threshold enough to recover, and from inspection of the per-round
checkpoints they do not move the underlying decision surface either.
This is consistent with the family-of-generators observation in
Grover \cite{zellers2019defending}: the best detector for an
out-of-family generator is often a detector that was \emph{also}
trained on that family.

\paragraph{Disclosure recommendation.}
We argue that any deployed biomedical AI-text detector should publish a
generator-family evasion matrix as part of its documentation, with at
least one in-family LLM, one out-of-family LLM, and one non-LLM
generator. Reporting only in-family AUC produces numbers that look
near-perfect and are operationally meaningless if the deployment
adversary uses any other generator.

\subsection{The C\,$=$\,D Pathology and Best-Val-AUC Selection}
\label{ssec:cd_path}

Conditions~C and~D produced byte-identical test metrics across all three
rounds. We diagnose this as an interaction between the detector's
ability and the silent default of best-validation-AUC checkpoint
selection.

The Condition~A baseline detector reaches validation AUC \(0.9990\) —
within \(\sim 0.001\) of the empirical ceiling on this corpus. In
Condition~D, each retraining round computes the validation AUC of
every epoch checkpoint and keeps the best. Because no retrained
checkpoint can exceed the baseline's \(0.9990\), the selector keeps
rolling back to the round-0 weights — i.e., to Condition~A's
detector. Condition~D's loop \emph{runs}, generates fakes, computes
hard subsets, calls the agent, and writes \(936\) new fakes to disk;
none of that work updates the deployed detector.

This is not a bug in our implementation. It is the expected behaviour
of best-val-AUC selection when the baseline is at ceiling. The
implication is that the metric one reports for ``the full pipeline''
is actually the metric of a detector that no rewrite or retrain ever
modified — an effect that, to our knowledge, is not flagged in the
adversarial-text-training literature.

\paragraph{Two ways out.}
\begin{enumerate}[leftmargin=*]
\item \emph{Adversarial-aware checkpoint selection}: keep the checkpoint
  with the lowest evasion on a held-out adversarial pool, not the
  highest validation AUC. This trades a small amount of clean-data
  performance for measurable adversarial signal.
\item \emph{Stronger generators}: an LLM that has been LoRA-fine-tuned
  to actively evade the detector (e.g., DPO with detector logits as
  reward) would push validation AUC below the baseline ceiling and
  give retraining a non-zero gradient. Our flat
  Condition~B curve suggests that without one of these two changes,
  the loop has nothing to learn from at this scale.
\end{enumerate}

\subsection{Cross-Domain Brittleness}
\label{ssec:domain_disc}

RAID's biomedical slice exposes a calibration failure that the
in-distribution metrics cleanly hide. AUC stays at \(0.82\), so the
detector retains ranking ability across domains. Accuracy collapses to
\(0.11\) and evasion shoots to \(0.983\), so its calibrated threshold
is meaningless on RAID's mixed-LLM fakes. The qualitative pattern is
the same brittleness reported by the RAID and MAGE papers; the
quantitative shape (huge AUC--accuracy gap) is a useful signal for
practitioners who only ever look at AUC.

For the BioShield deployment story, this means we cannot honestly
position the detector as a universal AI-text detector. We position it
instead as a \emph{biomedical specialist} trained against the
BioMistral family, with explicit cross-domain failure documented up
front. Reviewers asking for cross-domain robustness should be pointed
to the multi-family training proposal in Section~\ref{sec:limitations}.

\subsection{Retraction of the 500-Scale ``Adversarial Training Works''
Curve}
\label{ssec:retraction}

A previous version of this work, run at \(500\)-sample scale
(\(250\) real \(+\) \(250\) BioMistral fake, fake-pool
\(n = 250\) per round, test \(n = 100\)), reported a per-round
evasion curve of \(0.74 \to 0.64 \to 0.54\) on Condition~B and
interpreted it as evidence that 3-round adversarial retraining works
against SeqGAN-style fakes. The narrative slide deck for that run
opened with this number.

At \(3{,}000\)-sample scale, with the same
random seed, hyperparameters, and pipeline code, that curve does not
reproduce. Per-round evasion is \(0.993, 0.980, 0.987\) — flat within
the \(\pm 0.062\) Wilson CI on \(n = 250\). We retract the prior
``adversarial training works'' claim and attribute the \(500\)-scale
trend to small-sample noise: the per-round evasion sample is binary on
\(n = 250\), giving an effective standard deviation of
\(\sigma \approx 0.03\), which can produce a \(\Delta = 0.20\) drift
across three rounds with non-trivial probability under the null.

This is, in our reading, the most important methodological lesson of
the project: \emph{adversarial training results at sub-thousand sample
scale should be treated as exploratory, not load-bearing}, and any
published evasion curve should disclose its sample size and the CI
implied by it.

\section{Limitations and Future Work}
\label{sec:limitations}

We catalogue limitations explicitly so reviewers and downstream users can
calibrate the results.

\subsection{Limitations}

\paragraph{Single random seed per run.}
All four conditions and the RAID probe were evaluated under
\texttt{base=42} with per-round \texttt{set\_seed(42 + r)}. The
\(\pm 0.029\) Wilson half-width on test \(n=300\) bounds within-run
variance, but does not bound run-to-run variance from re-seeding.
Reported numbers should be read as point estimates; multi-seed
replication is required for any confidence-interval claim across
runs.

\paragraph{Three-thousand-sample scale.}
Test \(n=300\) produces a \(\pm 0.029\) CI half-width — wide enough that
the C\,$=$\,A and D\,$=$\,A equivalences are formally ``not significantly
different,'' rather than proven equal. A paper-grade replication should
push to \(n=2{,}000\) (Wilson half-width \(\pm 0.011\)) before treating
the C\,$=$\,D pathology claim as a tight finding.

\paragraph{Single biomedical domain.}
The detector is trained on PubMed and one biomedical-specialised LLM. We
do not test on clinical notes (MIMIC-III, MIMIC-IV) or on biomedical
preprints (bioRxiv, medRxiv). The cross-domain failure on RAID is
suggestive but RAID itself is heterogeneous; a within-biomedical
domain-shift study (PubMed-trained, MIMIC-tested) would tighten the
domain-shift claim.

\paragraph{Zero-shot generator and agent.}
Both BioMistral-7B (the generator) and Qwen2.5-7B (the agent) are used
zero-shot. A LoRA-fine-tuned version of either, trained against the
detector's logits, would likely change the picture. We treat this as
the most important pre-paper engineering investment
(Section~\ref{ssec:future}).

\paragraph{Best-validation-AUC checkpoint selection.}
The C\,$=$\,D pathology is downstream of this design choice. We have not
yet swept alternative selectors (best fake-pool evasion, best worst-case
evasion across pools, average over the last $k$ epochs). A small grid
sweep over selectors would convert our pathology observation into a
quantitative recommendation.

\paragraph{No human evaluation.}
We do not include a blinded human reader study in this draft. Reviewers
will reasonably want one; we plan a 50-sample blinded protocol
(Section~\ref{ssec:future}).

\paragraph{M4 cross-domain benchmark unavailable.}
The \texttt{Chardonneret/M4} dataset \cite{wang2024m4} was removed from
the Hugging~Face Hub between project start and this run. We log a
graceful skip in \texttt{condition\_D\_m4\_metrics.json} and propose
\texttt{Hello-SimpleAI/HC3} \cite{guo2023hc3} or
\texttt{yaful/MAGE} \cite{li2024mage} as replacements for the paper
revision.

\paragraph{Transfer attack queued, not run.}
A Llama-3.2-3B transfer attack is configured (model downloaded,
inference verified) but the \(500\)-sample fake corpus is not yet
generated. The completed family-cross-check matrix is the most
significant missing experiment.

\subsection{Future Work}
\label{ssec:future}

\paragraph{Pre-paper, this quarter.}
\begin{enumerate}[leftmargin=*]
\item \emph{Multi-seed replication.} Repeat the
  \(3{,}000\)-sample run under three additional seeds
  (\texttt{base} $\in \{43, 44, 45\}$) and report mean $\pm$ s.d.\
  rather than point estimates.
\item \emph{LoRA-hardened BioMistral.} Fine-tune BioMistral with DPO
  using the detector's real-probability score as the negative reward.
  Expected to break the validation-AUC ceiling and reactivate the
  retraining gradient.
\item \emph{Adversarial-aware checkpoint selector.} Replace
  best-val-AUC with best-fake-pool-evasion; re-run Conditions~B and~D
  end to end; report whether the C\,$=$\,D pathology dissolves.
\item \emph{Llama-3.2-3B transfer attack.} Generate \(500\) Llama
  fakes, score with Condition~A's detector, report a \(2\times 2\)
  family-cross-check matrix (BioMistral, Llama, SeqGAN, RAID).
\item \emph{Blinded human evaluation.} \(50\) abstracts (\(25\)
  real + \(25\) BioMistral fake), \(3\) M.Tech-level domain readers,
  per-reader and aggregated accuracy with Cohen's $\kappa$.
\end{enumerate}

\paragraph{Paper Tier~1 (peer-review defensible).}
Increase scale to \(2{,}000\) real + \(2{,}000\) fake, \(3\) rounds
$\times$ \(500\) per pool, test \(n=400\) (\(\pm 2.5\%\) CI). Wall-clock
budget on M3~Max: $\sim$\(16\) hours (one overnight). With batched
BioMistral generation (\texttt{batch\_size=4} in
\texttt{train\_generator.py}) we estimate
\(10\!-\!12\) hours.

\paragraph{Paper Tier~2 (reviewer-proof).}
Scale to \(10{,}000\) real + \(10{,}000\) fake, \(5\) rounds $\times$
\(1{,}000\) per pool, test \(n=2{,}000\) (\(\pm 1.1\%\) CI), full
family-cross-check matrix and human eval. Estimated wall-clock with
batching: $\sim$\(40\!-\!65\) hours.

\paragraph{Beyond the current pipeline.}
Two longer-horizon directions are worth flagging. (i)~A multi-domain
detector trained jointly on PubMed, MIMIC, and bioRxiv with
domain-adversarial pretraining \cite{ganin2016dann} should help with
cross-domain calibration. (ii)~A multi-family training corpus including
SeqGAN-style fakes, BioMistral fakes, and at least one general-purpose
LLM family (Llama, GPT) would convert the family-disclosure
recommendation in Section~\ref{ssec:family_disc} from a measurement
into a deployable defence.

\section{Conclusion}
\label{sec:conclusion}

We set out to test whether a 3-round adversarial retraining loop, with
or without an LLM rewriting agent, can harden a BiomedBERT detector
against AI-generated biomedical abstracts. On a controlled
\(3{,}000\)-sample corpus, evaluated under fixed seeds and
hyperparameters on a single Apple M3~Max laptop in
\(10\,\mathrm{h}\,01\,\mathrm{m}\) of wall-clock, the answer is no — and
the negative result is informative.

Three findings stand out. First, the choice of \emph{generator family}
explains roughly \(140\times\) more variance in evasion rate than the
choice of retraining schedule or rewriting agent: BioMistral fakes evade
at \(0.7\%\), SeqGAN fakes at \(98.7\%\), RAID's mixed LLMs at
\(98.3\%\), all on the same detector. Second, three rounds of
adversarial retraining produced no measurable improvement on either
in-family (Condition~C, D) or out-of-family (Condition~B) fakes; the
common best-validation-AUC checkpoint selector silently rolled
Conditions~C and~D back to the baseline weights, which is the source of
their byte-identical metrics. Third, a previous \(500\)-sample evasion
curve of \(0.74 \to 0.54\) in Condition~B did not reproduce at
\(3{,}000\) scale and is best explained as small-sample noise within the
\(\pm 0.062\) Wilson CI on \(n = 250\).

The actionable takeaways for downstream practitioners are concrete: any
deployed AI-text detector should publish a
generator-family evasion matrix as part of its documentation; the
default best-validation-AUC checkpoint selector should be replaced by an
adversarial-aware selector when the baseline is at ceiling; and
sub-thousand-sample evasion claims should be treated as exploratory
until they are replicated at a scale where Wilson CIs do not absorb the
effect being claimed. We position BioShield as a biomedical specialist
detector with documented cross-domain failure, release the full code,
data, metrics, and one-command rerun script at
\url{https://github.com/Rayudualapati-25/BioShield}, and treat the
multi-seed replication, LoRA-hardened generator, and family-cross-check
matrix as the highest-priority follow-ups for the paper revision.


\section*{Reproducibility Statement}
All code, configuration files, per-round metric JSONs, plots, and the
\(3{,}000\)-sample data manifest are released at
\url{https://github.com/Rayudualapati-25/BioShield}. The complete
\(10\,\mathrm{h}\,01\,\mathrm{m}\) pipeline reproduces with
\texttt{bash scripts/run\_presentation.sh} on a single Apple~M3~Max~64\,GB
machine. Random seeds (\texttt{base=42}, per-round
\texttt{set\_seed(42+r)}), bf16 dtype, and the
\texttt{microsoft/BiomedNLP-BiomedBERT-large} checkpoint are fixed across
conditions; the only variable is training data.

\section*{Acknowledgment}
The author thanks the maintainers of the BioMistral, BiomedBERT, Qwen, and
SeqGAN projects for releasing their models and code, and the Hugging~Face Hub
team for streaming PubMed access. The RAID benchmark
\cite{dugan2024raid} provided the cross-domain probe that anchors our
domain-shift analysis.

\bibliographystyle{IEEEtran}
\bibliography{refs}

\end{document}
